\chapter{Séries de potências}\label{ch:b2:powerseries}

As séries de potências são as séries de funções mais bem-comportadas da
matemática: dentro de seu disco de convergência elas convergem
\omterm{def:b2:funcseq:series}{normalmente} nos \omterm{def:b2:metric:compact}{compactos}, podem ser derivadas e integradas termo a
termo sem pensar duas vezes, e suas somas --- as \emph{\omterm{def:b2:powerseries:analytic}{funções
analíticas}} --- são determinadas por seus coeficientes. Este capítulo
demonstra esse pacote inteiro e recupera, honestamente, toda série de
Taylor do volume do primeiro ano de graduação; as \omterm{ex:b2:powerseries:fibonacci}{funções geradoras} o fecham com
dividendos algébricos.

\section{Raio de convergência}

\begin{lemma}[Abel]\label{lem:b2:powerseries:abel}
Se a sequência $(a_n z_0^n)$ é limitada para algum $z_0 \neq 0$,
então $\sum a_n z^n$ converge \omterm{def:b2:series:def}{absolutamente} para todo $\abs z < \abs
{z_0}$, e \omterm{def:b2:funcseq:series}{normalmente} em todo disco $\abs z \leq r < \abs{z_0}$.
\end{lemma}

\begin{proof}
Com $\abs{a_n z_0^n} \leq M$ e $\abs z \leq r$:
\[
\abs{a_n z^n} = \abs{a_n z_0^n}\,\Bigl|\frac{z}{z_0}\Bigr|^n
\leq M\Bigl(\frac{r}{\abs{z_0}}\Bigr)^{\!n},
\]
uma cota geométrica convergente, uniforme no disco.
\end{proof}

\begin{definition}[Raio de convergência]\label{def:b2:powerseries:radius}
O \emph{raio de convergência}\index{raio de convergência} de
$\sum a_n z^n$ é
\[
R = \sup\{r \geq 0 : (a_n r^n) \text{ limitada}\} \in
\intcc{0}{+\infty} .
\]
Pelo~\cref{lem:b2:powerseries:abel}: convergência absoluta para $\abs
z < R$ (normal em subdiscos \omterm{def:b2:metric:compact}{compactos}), divergência --- com termos até
ilimitados --- para $\abs z > R$. No círculo de fronteira, tudo pode
acontecer (\cref{exo:b2:powerseries:2}). Na prática, $R$ é calculado
pelo critério da razão de d'Alembert em $\abs{a_n}\abs z^n$ ou por comparação.
\end{definition}

\begin{example}[Um raio sem critério da razão]\label{ex:b2:powerseries:sinn}
Qual é o raio de $\sum \sin(n)\,z^n$? A razão
$\abs{\sin(n+1)/\sin n}$ não tem limite, mas a definição funciona
diretamente. \emph{$R \geq 1$:} $\abs{\sin n} \leq 1$, logo $(\sin
n\cdot r^n)$ é limitada para todo $r < 1$ --- de fato, para $r =
1$. \emph{$R \leq 1$:} basta que $\sin n \not\to 0$.
Suponha $\sin n \to 0$; a fórmula de adição
\[
\sin(n+1) = \sin n\cos 1 + \cos n\sin 1
\]
forçaria $\cos n \to 0$ (resolva para $\cos n$: $\sin 1 \neq
0$), contradizendo $\sin^2 n + \cos^2 n = 1$. Assim os termos
$\sin(n)\,1^n$ não tendem a $0$: a série diverge em $z =
1$, e $R \leq 1$. Conclusão: $R = 1$. Lição final: o
raio é um enunciado sobre $\abs{a_n}r^n$ ser limitada ---
nenhum limite de razões é jamais exigido, e argumentos de limitação
resolvem casos que o critério da razão não consegue tocar (compare os
coeficientes oscilantes do~\cref{exo:b2:powerseries:1}).
\end{example}

\begin{proposition}[Operações]\label{prop:b2:powerseries:operations}
Sejam $\sum a_nz^n$ e $\sum b_nz^n$ de raios $R_a, R_b$. Então, para
$\abs z < \min(R_a, R_b)$:
\[
\sum (a_n + b_n)z^n = \sum a_nz^n + \sum b_nz^n,
\qquad
\Bigl(\sum a_nz^n\Bigr)\Bigl(\sum b_nz^n\Bigr) = \sum c_n z^n,
\quad c_n = \sum_{k=0}^{n} a_kb_{n-k},
\]
tendo as duas séries raio $\geq \min(R_a, R_b)$. (O produto é o
\omterm{thm:b2:series:fubini}{produto de Cauchy}, legítimo pela convergência absoluta e pelo
\cref{thm:b2:series:fubini}.)
\end{proposition}

\begin{proof}
A fórmula da soma é a linearidade das séries convergentes, e
$(a_n + b_n)r^n$ é limitada sempre que $a_nr^n$ e $b_nr^n$
o são: raio $\geq \min(R_a, R_b)$. Para o produto, fixe $\abs z
< \min(R_a, R_b)$: as duas séries convergem \emph{\omterm{def:b2:series:def}{absolutamente}}
aí (\cref{lem:b2:powerseries:abel}), de modo que a família duplamente indexada
$(a_kz^k\,b_lz^l)_{k,l}$ é \omterm{def:b2:series:summable}{somável}, e o
\cref{thm:b2:series:fubini} permite qualquer agrupamento. Agrupando por $k
+ l = n$:
\[
\Bigl(\sum_k a_kz^k\Bigr)\Bigl(\sum_l b_lz^l\Bigr)
= \sum_{n\geq0}\Bigl(\sum_{k+l=n}a_kb_l\Bigr)z^n
= \sum_{n\geq0}c_nz^n ,
\]
\omterm{def:b2:series:def}{absolutamente} convergente para todo tal $z$: a série produto
tem raio $\geq \min(R_a, R_b)$ também.
\end{proof}

\begin{example}[Um quadrado de Cauchy, conferido]\label{ex:b2:powerseries:cauchysquare}
Eleve ao quadrado a série geométrica: para $\abs x < 1$, o coeficiente
de $x^n$ em $\bigl(\sum x^k\bigr)^2$ é $c_n = \sum_{k+l=n}
1\cdot1 = n + 1$, logo
\[
\frac{1}{(1-x)^2} = \sum_{n\geq0}(n+1)\,x^n .
\]
Confira por derivação termo a termo
(\cref{thm:b2:powerseries:calculus} adiante): derivando
$\frac{1}{1-x} = \sum x^n$ obtém-se $\frac{1}{(1-x)^2} = \sum
nx^{n-1} = \sum(n+1)x^n$ --- a mesma série por dois
mecanismos independentes. Lição final: quando uma identidade entre
coeficientes parece misteriosa, um desses dois motores (convolução ou
derivação) costuma produzi-la em uma linha; a questão do problema de
fim de semana sobre $\sum\binom{2k}k\binom{2n-2k}{n-k} = 4^n$
roda o motor de convolução a plena potência.
\end{example}

\begin{example}[Multiplicar por $\frac{1}{1-x}$ soma os coeficientes]\label{ex:b2:powerseries:partialsums}
Um \omterm{thm:b2:series:fubini}{produto de Cauchy} contra a série geométrica tem um
significado memorável: para qualquer $\sum a_nx^n$ de raio $R > 0$ e $\abs x <
\min(R, 1)$,
\[
\frac{1}{1-x}\sum_{n\geq0}a_nx^n
= \sum_{n\geq0}\Bigl(\sum_{k=0}^{n}a_k\Bigr)x^n :
\]
multiplicar por $\frac{1}{1-x}$ substitui os coeficientes por
suas somas parciais (convolução com a sequência de uns).
Exemplo: $\dfrac{\eu^x}{1-x} = \sum_n s_n x^n$ com $s_n =
\sum_{k\leq n}\frac{1}{k!}$, as somas parciais de $\eu$ ---
compare com o~\cref{exo:b2:powerseries:11}, em que o mesmo produto
com $\eu^{-x}$ codifica as contagens de desarranjos. Lição
final: operações sobre séries de potências são operações sobre
sequências de coeficientes disfarçadas (multiplicar por $\frac1{1-x}$:
somar; multiplicar por $x$: deslocar; derivar: multiplicar por $n$
e deslocar) --- um dicionário que o capítulo de \omterm{ex:b2:powerseries:fibonacci}{funções geradoras}
lerá com fluência.
\end{example}

\section{Regularidade da soma}

\begin{theorem}[Cálculo termo a termo]\label{thm:b2:powerseries:calculus}
Seja $f(x) = \sum_{n\geq0} a_n x^n$ de raio $R > 0$ (variável
real $x \in \intoo{-R}{R}$).
\begin{enumerate}
  \item A série derivada $\sum n\,a_n x^{n-1}$ tem o \emph{mesmo}
        raio $R$, e $f$ é $C^1$ com $f'(x) = \sum_{n \geq 1}
        n a_n x^{n-1}$. Iterando, $f$ é $C^\infty$ e
        \[
        a_n = \frac{f^{(n)}(0)}{n!} :
        \]
        os coeficientes de uma série de potências são únicos (duas séries
        com a mesma soma perto de $0$ têm coeficientes iguais).
  \item Primitiva termo a termo: $\sum \frac{a_n}{n+1}x^{n+1}$ tem
        raio $R$ e derivada $f$.
\end{enumerate}
\end{theorem}

\begin{proof}
\emph{Mesmo raio:} se $(a_nr^n)$ é limitada e $r' < r$, então
$n\abs{a_n} r'^{\,n-1} = \frac{n}{r'}\abs{a_nr^n}
\bigl(\frac{r'}{r}\bigr)^n$ é limitada (de fato $\to 0$: a geométrica
vence $n$), logo $R' \geq R$; reciprocamente, $\abs{a_n x^n} \leq \abs x
\cdot n\abs{a_n}\abs x^{n-1}$ dá $R \geq R'$.

\emph{Derivação:} em $\intcc{-r}{r}$, $r < R$, a série
derivada converge \omterm{def:b2:funcseq:series}{normalmente} ($n\abs{a_n}r^{n-1}$ \omterm{def:b2:series:summable}{somável} pelo
cálculo do raio); a original converge em $x = 0$: o
teorema de derivação para séries
(\cref{thm:b2:funcseq:seriestransfer}) se aplica em todo segmento
desse tipo, logo em $\intoo{-R}{R}$. Iterando $k$ vezes e
avaliando em $0$: explicitamente, a $k$-ésima série derivada é
\[
f^{(k)}(x) = \sum_{n\geq k} n(n-1)\cdots(n-k+1)\,a_n\,x^{n-k},
\]
e em $x = 0$ todo termo com $n > k$ se anula, restando apenas
o termo constante $k(k-1)\cdots1\cdot a_k$: $f^{(k)}(0) =
k!\,a_k$. A unicidade dos coeficientes segue: duas séries de potências
com a mesma soma perto de $0$ têm as mesmas derivadas em $0$,
logo os mesmos $a_k$. Primitivas: mesmo raio pelo mesmo
cálculo, e derive termo a termo de volta.
\end{proof}

\begin{example}[Avaliar uma série num ponto]\label{ex:b2:powerseries:evaluate}
Quanto vale $\sum_{n\geq1}\dfrac{n^2}{2^n}$? É a soma
$\sum n^2x^n$ do~\cref{exo:b2:powerseries:3} avaliada
\emph{dentro} do disco, em $x = \frac12 < 1 = R$, onde toda
manipulação usada para deduzir a forma fechada era legítima:
\[
\sum_{n\geq1} n^2x^n = \frac{x(1+x)}{(1-x)^3}
\quad\Longrightarrow\quad
\sum_{n\geq1}\frac{n^2}{2^n}
= \frac{\frac12\cdot\frac32}{(\frac12)^3}
= \frac{3/4}{1/8} = 6 .
\]
Mesmo motor, outros botões: $x = \frac13$ dá
$\sum\frac{n^2}{3^n} = \frac{\frac13\cdot\frac43}{(2/3)^3} =
\frac32$. Lição final: uma identidade de séries de potências é uma
máquina, e não uma única fórmula --- uma dedução precifica de uma vez
toda série numérica $\sum n^2q^n$, para todo $\abs q < 1$;
é assim que o capítulo de \omterm{ex:b2:powerseries:fibonacci}{funções geradoras} vai calcular
esperanças e variâncias no atacado.
\end{example}

\begin{example}[Os clássicos, desta vez honestamente]\label{ex:b2:powerseries:classics}
$\displaystyle\frac{1}{1 - x} = \sum x^n$ ($R = 1$); integrando
termo a termo (\cref{thm:b2:powerseries:calculus}~(2)):
\[
-\ln(1 - x) = \sum_{n\geq1} \frac{x^n}{n},
\qquad
\arctan x = \sum_{n \geq 0} \frac{(-1)^n x^{2n+1}}{2n+1}
\quad (\abs x < 1),
\]
a segunda em dois passos: substitua $-x^2$ na série
geométrica para obter $\frac{1}{1+x^2} = \sum(-1)^nx^{2n}$ (raio
$1$, pois $\abs{x^2} < 1 \iff \abs x < 1$) e depois tome a
primitiva termo a termo que se anula em $0$; os dois lados são
primitivas da mesma função com o mesmo valor em $0$,
logo iguais em $\intoo{-1}{1}$. E $\exp$:
a série $E(x) = \sum \frac{x^n}{n!}$ ($R = \infty$) satisfaz
$E' = E$, $E(0) = 1$ por derivação termo a termo, logo $E = \exp$ pela
unicidade do primeiro ano. Todo ``desenvolvimento padrão'' do volume do primeiro
ano de graduação é agora um teorema sobre sua série de potências \omterm{def:b2:metric:complete}{completa}.
\end{example}

\begin{example}[Um logaritmo calculado de dentro do disco]\label{ex:b2:powerseries:lntwo}
Avaliando $-\ln(1-x) = \sum\frac{x^n}{n}$ no ponto
interior $x = \frac12$:
\[
\sum_{n\geq1}\frac{1}{n\,2^n} = \ln 2 ,
\]
uma representação de rápida convergência de $\ln 2$ (dez termos já
dão $0.69306\ldots$ contra $\ln 2 = 0.69314\ldots$), bem
melhor do que a série alternada $1 - \frac12 + \frac13 -
\dots$, disponível apenas na fronteira. Lição final:
sempre que uma constante for atingível tanto na borda quanto estritamente
dentro do disco, o interior vence numericamente --- decaimento
geométrico contra decaimento harmônico.
\end{example}

\begin{example}[A derivação preserva o raio, não a fronteira]\label{ex:b2:powerseries:boundaryloss}
A série $\sum_{n\geq1}\frac{x^n}{n^2}$ tem raio $1$ e
converge nas \emph{duas} extremidades ($\sum\frac1{n^2}$ e sua
gêmea alternada). Sua série derivada,
\[
\sum_{n\geq1}\frac{x^{n-1}}{n} ,
\]
tem o mesmo raio $1$ --- como
o~\cref{thm:b2:powerseries:calculus} garante --- mas agora
diverge em $x = 1$ (série harmônica) enquanto ainda converge
em $x = -1$ (alternada). Mais uma derivação dá
$\sum_{n\geq2}\frac{n-1}{n}x^{n-2}$, divergente nas duas pontas
(os termos não tendem a $0$). Lição final: cada
derivação multiplica os coeficientes por $n$, o que
nunca move o raio (a geométrica vence a polinomial) mas devora
uma ordem de decaimento na fronteira; o cálculo termo a termo é um
esporte de interior, e o que quer que aconteça na borda tem de ser
reexaminado --- a teoria de Abel--Tauber do problema de fim de semana é
exatamente esse reexame.
\end{example}

\begin{example}[Separar uma série por restos --- trabalhada até o fim]\label{ex:b2:powerseries:quarter}
Calcule $f(x) = \sum_{n\geq0} \dfrac{x^{4n}}{(4n)!}$ em forma
fechada. Tanto $\cosh x = \sum \frac{x^{2m}}{(2m)!}$ quanto $\cos x =
\sum \frac{(-1)^m x^{2m}}{(2m)!}$ têm raio $\infty$, de modo que sua
média pode ser calculada termo a termo:
\[
\frac{\cosh x + \cos x}{2}
= \sum_{m\geq0}\frac{1 + (-1)^m}{2}\,\frac{x^{2m}}{(2m)!}
= \sum_{m \text{ par}}\frac{x^{2m}}{(2m)!}
= \sum_{n\geq0}\frac{x^{4n}}{(4n)!} = f(x) .
\]
O filtro $\frac{1+(-1)^m}{2}$ retém exatamente os $m$ pares: esse
é o avatar real do filtro das raízes da unidade (a versão
complexa, com $\iu^n$, extrai restos módulo $4$ de um só
golpe). Verificação final: $f$ resolve $f'''' = f$ com $f(0) = 1$,
$f'(0) = f''(0) = f'''(0) = 0$ --- derive a série quatro
vezes (\cref{thm:b2:powerseries:calculus}) e veja-a
reproduzir-se; $\frac{\cosh + \cos}{2}$ satisfaz os mesmos
dados.
\end{example}

\begin{definition}[Funções analíticas]\label{def:b2:powerseries:analytic}
$f$ é \emph{analítica}\index{função analítica} em $x_0$ quando é
a soma de uma série de potências em $(x - x_0)$ numa vizinhança; num
intervalo, quando o é em todo ponto. As somas de séries de potências são analíticas
dentro de seu disco (reordenação do desenvolvimento --- admitida
neste nível para o recentramento, sendo o caso $x_0 = 0$ o
\cref{thm:b2:powerseries:calculus}). Analítica implica $C^\infty$;
a recíproca \emph{falha}: a função achatada $\eu^{-1/x^2}$
(\cref{exo:b2:powerseries:7}).
\end{definition}

\begin{example}[Recentrar, e o raio como distância]\label{ex:b2:powerseries:recenter}
Desenvolva $f(x) = \frac{1}{1-x}$ em torno de $x_0 = \frac12$: escrevendo
$x = \frac12 + h$,
\[
\frac{1}{1 - x} = \frac{1}{\frac12 - h}
= \frac{2}{1 - 2h}
= \sum_{n\geq0} 2^{n+1}\,h^n
= \sum_{n\geq0} 2^{n+1}\Bigl(x - \frac12\Bigr)^{\!n},
\]
válido para $\abs{2h} < 1$, isto é, $\abs{x - \frac12} < \frac12$.
O novo raio é exatamente a distância do novo centro à
singularidade $x = 1$: recentrar encolhe (ou aumenta) o
disco para que ele caiba até o obstáculo mais próximo. Lição final: essa é
a figura por trás da definição de analiticidade --- uma
função, muitas séries de potências locais, cada uma vivendo no maior
disco que evita o problema; o volume do terceiro ano de graduação transforma a
heurística ``raio $=$ distância à singularidade complexa mais
próxima'' num teorema.
\end{example}

\begin{remark}[Armadilhas comuns]
\emph{(i) O critério da razão é suficiente, não necessário:} quando
$\abs{a_{n+1}/a_n}$ não tem limite
(\cref{ex:b2:powerseries:sinn}, \cref{exo:b2:powerseries:1}),
volte à definição: $R = \sup\{r : (a_nr^n)$
limitada$\}$. \emph{(ii) Nada atravessa a fronteira de graça:}
a derivação e a integração termo a termo são teoremas
\emph{dentro} do disco \omterm{def:b2:metric:topology}{aberto}; em $\abs x = R$ cada série tem de
ser reexaminada (esse é todo o assunto do problema de fim
de semana). \emph{(iii) Raio de uma soma:} $\min(R_a, R_b)$ é
apenas uma cota inferior --- cancelamentos podem aumentá-lo ($a_n = 1,
b_n = -1$: soma identicamente $0$, raio $\infty$). \emph{(iv)
$C^\infty$ não é \omterm{def:b2:powerseries:analytic}{analítica}:} uma série de Taylor convergente pode
convergir para a função \emph{errada}
(\cref{exo:b2:powerseries:7}); antes de escrever $f(x) = \sum
\frac{f^{(n)}(0)}{n!}x^n$, demonstre-o --- por uma EDO
(\cref{met:b2:powerseries:ode}), uma estimativa do resto ou uma
fórmula integral.
\end{remark}

\begin{remark}[Onde isso é usado]
As séries de potências são o cavalo de batalha de três capítulos posteriores: o
capítulo de equações diferenciais resolve EDOs lineares injetando
$\sum a_nx^n$ (a caixa de método abaixo, industrializada); o
capítulo de \omterm{ex:b2:powerseries:fibonacci}{funções geradoras} converte identidades sobre
probabilidades em identidades sobre séries de potências e de volta; e
o volume do terceiro ano de graduação oficializa a variável complexa, onde a
analiticidade se torna equivalente à derivabilidade complexa e
o ``recentramento admitido'' acima ganha sua demonstração honesta. O
problema de fim de semana explora o único lugar em que os teoremas deste capítulo
ficam calados: a própria fronteira $\abs x = R$.
\end{remark}

\begin{method}[Desenvolver por meio de uma equação diferencial]\label{met:b2:powerseries:ode}
Para desenvolver uma função $f$ em série de potências: ache uma EDO linear com
coeficientes polinomiais satisfeita por $f$; injete $\sum a_nx^n$;
identifique os coeficientes para obter uma recorrência para $(a_n)$; resolva e
confira o raio e as condições iniciais. Exemplo --- a série
binomial: $f(x) = (1+x)^\alpha$ satisfaz $(1+x)f' = \alpha f$, $f(0)
= 1$; injetando obtém-se $(n+1)a_{n+1} = (\alpha - n)a_n$, logo $a_n =
\binom{\alpha}{n}$, raio $1$ (critério da razão) e a soma, satisfazendo
a mesma EDO com o mesmo valor inicial, é igual a $(1 + x)^\alpha$
pelo teorema de unicidade para EDOs lineares (volume do primeiro ano de graduação).
\end{method}

\begin{example}[O método numa equação com termo forçante]\label{ex:b2:powerseries:odeforced}
Resolva $y' = y + x$, $y(0) = 0$, por séries de potências. Injetando $y =
\sum a_nx^n$:
\[
\sum_{n\geq0}(n+1)a_{n+1}x^n
= \sum_{n\geq0}a_nx^n + x ,
\]
e identificando coeficientes: $a_1 = a_0 = 0$, $2a_2 = a_1 + 1
= 1$ e $(n+1)a_{n+1} = a_n$ para $n \geq 2$. Logo $a_2 =
\frac{1}{2!}$ e, por indução, $a_n = \frac{1}{n!}$ para todo
$n \geq 2$: raio $\infty$, e
\[
y(x) = \sum_{n\geq2}\frac{x^n}{n!} = \eu^x - 1 - x .
\]
Verificação: $y' = \eu^x - 1 = y + x$ e $y(0) = 0$. Lição
final: a recorrência \emph{é} a equação, coeficiente a
coeficiente; o termo forçante só perturba finitos coeficientes
iniciais, após os quais o padrão homogêneo assume o comando
--- uma sombra discreta de ``solução particular mais
solução homogênea''.
\end{example}

\section{Funções geradoras}

\begin{example}[Fibonacci]\label{ex:b2:powerseries:fibonacci}
Seja $F(x) = \sum_{n\geq0} F_n x^n$ (números de Fibonacci, $F_0 = 0$,
$F_1 = 1$). A recorrência $F_{n+2} = F_{n+1} + F_n$ traduz-se,
multiplicando por $x^{n+2}$ e somando, em
\[
F(x) - x = x\,F(x) + x^2 F(x)
\quad\Longrightarrow\quad
F(x) = \frac{x}{1 - x - x^2} ,
\]
válido onde a série converge. O raio é
$\frac{1}{\varphi}$: de $F_n \sim
\frac{\varphi^n}{\sqrt5}$ (Binet, exemplo seguinte --- ou a indução
bruta $F_n \leq 2^n$ mais a recorrência), o critério da razão
dá
\[
\frac{F_{n+1}\abs x^{n+1}}{F_n\abs x^n}
\longrightarrow \varphi\abs x ,
\qquad\text{convergência se e somente se } \abs x < \frac1\varphi
\approx 0.618 .
\] Frações parciais em
$\frac{x}{1 - x - x^2}$ e a série geométrica redemonstram a fórmula
de Binet --- as \omterm{ex:b2:powerseries:fibonacci}{funções geradoras} industrializam as recorrências
lineares.\index{função geradora}
\end{example}

\begin{example}[A fórmula de Binet, executada]\label{ex:b2:powerseries:binet}
Sejam $\varphi = \frac{1+\sqrt5}{2}$ e $\psi =
\frac{1-\sqrt5}{2}$, as raízes de $X^2 = X + 1$; como $\varphi
+ \psi = 1$ e $\varphi\psi = -1$,
\[
1 - x - x^2 = (1 - \varphi x)(1 - \psi x) .
\]
Frações parciais: procurando $\frac{x}{(1-\varphi x)(1-\psi x)} =
\frac{A}{1 - \varphi x} + \frac{B}{1 - \psi x}$, o termo
constante dá $A + B = 0$ e o coeficiente de $x$ dá $-A\psi - B\varphi
= 1$, logo $A(\varphi - \psi) = 1$: $A = \frac{1}{\sqrt5} = -B$.
Duas séries geométricas depois,
\[
F(x) = \frac{1}{\sqrt5}\sum_{n\geq0}
\bigl(\varphi^n - \psi^n\bigr)x^n
\quad\Longrightarrow\quad
F_n = \frac{\varphi^n - \psi^n}{\sqrt5}
\]
pela unicidade dos coeficientes
(\cref{thm:b2:powerseries:calculus}). Como $\abs\psi < 1$, o
termo $\frac{\psi^n}{\sqrt5}$ tem valor absoluto $< \frac12$:
$F_n$ é o \emph{inteiro mais próximo} de
$\frac{\varphi^n}{\sqrt5}$. Lição final: o raio
$\frac1\varphi$ de $F$ é o inverso da raiz dominante
--- o crescimento dos coeficientes e o \omterm{def:b2:powerseries:radius}{raio de convergência} são a
mesma informação lida em sentidos opostos.
\end{example}

\begin{example}[Números de Catalan]\label{ex:b2:powerseries:catalan}
Os \omterm{ex:b2:powerseries:catalan}{números de Catalan} $C_n$ (número de triangulações, de parentizações,
de caminhos de Dyck, \dots) satisfazem $C_0 = 1$ e $C_{n+1} = \sum_{k=0}^n
C_kC_{n-k}$. A \omterm{ex:b2:powerseries:fibonacci}{função geradora} $C(x) = \sum C_nx^n$ satisfaz então (\omterm{thm:b2:series:fubini}{produto de Cauchy}!)
\[
C(x) = 1 + x\,C(x)^2
\quad\Longrightarrow\quad
C(x) = \frac{1 - \sqrt{1 - 4x}}{2x} ,
\]
escolhendo a raiz com $C(0) = 1$: resolver a quadrática
$xC^2 - C + 1 = 0$ dá os dois candidatos $\frac{1 \pm
\sqrt{1-4x}}{2x}$ e, quando $x \to 0$, a raiz ``$+$'' explode
como $\frac1x$ enquanto a raiz ``$-$'' tende a $1$ (desenvolva
$\sqrt{1-4x} = 1 - 2x + O(x^2)$) --- só o sinal de menos pode
carregar uma série de potências com $C_0 = 1$. Desenvolver $\sqrt{1 - 4x}$ pela
série binomial dá a forma fechada
\[
C_n = \frac{1}{n+1}\binom{2n}{n} ,
\]
executada no~\cref{exo:b2:powerseries:8}.\index{números de Catalan}
\end{example}

\begin{remark}[Séries formais contra séries convergentes]
Todo cálculo com \omterm{ex:b2:powerseries:fibonacci}{funções geradoras} acima termina invocando
a unicidade dos coeficientes, e esse teorema mora
\emph{dentro} de um disco de raio positivo: antes de ``ler''
$F_n$ ou $C_n$, é preciso saber que $R > 0$. Uma cota a priori bruta
basta --- $F_n \leq 2^n$ (indução imediata) dá $R \geq
\frac12$ para Fibonacci; $C_n \leq 4^n$ (cada \omterm{ex:b2:powerseries:catalan}{número de Catalan}
conta subconjuntos de caminhos) dá $R \geq \frac14$. Cuidado com a
ponta degenerada da escala: $\sum n!\,x^n$ tem raio $0$, e
manipulá-la como função não faz sentido --- identidades
envolvendo tais séries pertencem ao cálculo \emph{formal} dos
coeficientes, um jogo puramente algébrico com suas próprias (outras)
regras. Neste nível: garanta sempre primeiro um raio positivo e
depois calcule livremente dentro dele.
\end{remark}

\begin{remark}[Perspectivas dentro deste volume]
As séries de potências são uma das duas grandes máquinas de expansão
do livro; a outra é a série de Fourier dos capítulos harmônicos,
e compará-las é instrutivo. Uma série de potências é
rígida: seus coeficientes são forçados ($a_n = f^{(n)}(0)/n!$),
sua convergência é implacável (normal dentro, impossível fora)
e sua soma é \omterm{def:b2:powerseries:analytic}{analítica} --- infinitamente rígida
(\cref{def:b2:powerseries:analytic}). Uma série de Fourier é
flexível: ela representa meros sinais suaves por partes, ao
preço de delicadas questões de convergência na fronteira da
suavidade. As duas teorias se encontram no problema de fim de semana
deste capítulo: as médias de Cesàro e a \omterm{thm:b2:series:abel}{transformação de Abel}, desenvolvidas aqui para
o círculo de fronteira, voltam no capítulo de Fourier como os núcleos de
Fejér e de Poisson. Enquanto isso, o capítulo de equações diferenciais
consome séries de potências diretamente ($\eu^{tA}$, soluções
em série), e o capítulo de \omterm{ex:b2:powerseries:fibonacci}{funções geradoras} transforma
o truque do~\cref{ex:b2:powerseries:fibonacci} num cálculo sistemático
para probabilidades.
\end{remark}

\section{Exercícios}

\begin{exercise}[$\star$]\label{exo:b2:powerseries:1}
\omterm{def:b2:powerseries:radius}{Raios de convergência}:
$\sum \dfrac{n^2}{2^n}z^n$; $\;\sum \dfrac{z^n}{\binom{2n}{n}}$;
$\;\sum z^{n!}$; $\;\sum \bigl(2 + (-1)^n\bigr)^n z^n$.
\end{exercise}

\begin{exercise}[$\star$]\label{exo:b2:powerseries:2}
Mostre que $\sum z^n$, $\sum \frac{z^n}{n}$ e $\sum \frac{z^n}{n^2}$
têm todas raio $1$ mas se comportam de modo diferente em $z = 1$ e $z = -1$:
divergência/divergência, divergência/convergência,
convergência/convergência.
\end{exercise}

\begin{exercise}[$\star$]\label{exo:b2:powerseries:3}
Calcule as somas, para $\abs x < 1$:
\[
\sum_{n\geq0} n x^n,
\qquad
\sum_{n\geq0} n^2 x^n,
\qquad
\sum_{n\geq0} \frac{x^{2n+1}}{2n+1} .
\]
\end{exercise}

\begin{exercise}[$\star\star$]\label{exo:b2:powerseries:4}
Desenvolva em série de potências em $0$, com o raio: $\dfrac{1}{(1-x)(2-x)}$
(frações parciais); $\;\ln(1 + x + x^2)$ \emph{(escreva $1 + x + x^2
= \frac{1 - x^3}{1 - x}$)}.
\end{exercise}

\begin{exercise}[$\star\star$]\label{exo:b2:powerseries:5}
Prove que $f(x) = \sum_{n\geq1} H_n x^n = -\dfrac{\ln(1 -
x)}{1 - x}$ para $\abs x < 1$, em que $H_n$ é o número harmônico
\emph{(\omterm{thm:b2:series:fubini}{produto de Cauchy} de $\sum x^n$ e $\sum \frac{x^n}{n}$)}.
\end{exercise}

\begin{exercise}[$\star\star$]\label{exo:b2:powerseries:6}
Resolva por \omterm{ex:b2:powerseries:fibonacci}{função geradora} a recorrência $u_0 = 1$, $u_{n+1} =
2u_n + n$: calcule $U(x) = \sum u_nx^n$ em forma fechada, decomponha
e leia $u_n = 2^{n+1} - n - 1$.
\end{exercise}

\begin{exercise}[$\star\star$]\label{exo:b2:powerseries:7}
Seja $f(x) = \eu^{-1/x^2}$ para $x \neq 0$, $f(0) = 0$. Prove que
$f$ é $C^\infty$ em $\R$ com $f^{(n)}(0) = 0$ para todo $n$
\emph{(mostre por indução que $f^{(n)}(x) = P_n\bigl(\frac1x\bigr)
\eu^{-1/x^2}$ para polinômios $P_n$, e use a comparação
de crescimento)}. Conclua que $f$ não é \omterm{def:b2:powerseries:analytic}{analítica} em $0$: sua série
de Taylor em $0$ converge --- para a função errada.
\end{exercise}

\begin{exercise}[$\star\star\star$]\label{exo:b2:powerseries:8}
Complete o~\cref{ex:b2:powerseries:catalan}: desenvolva $\sqrt{1 - 4x}$
com a série binomial, mostrando que
\[
\binom{1/2}{n+1}(-4)^{n+1} = -\frac{2}{n+1}\binom{2n}{n},
\]
e deduza $C_n = \frac{1}{n+1}\binom{2n}{n}$; determine o \omterm{def:b2:powerseries:radius}{raio
de convergência} de $C(x)$ e a assintótica de $C_n$ via
Stirling.
\end{exercise}

\begin{exercise}[$\star\star\star$]\label{exo:b2:powerseries:9}
(Teorema do limite radial de Abel, caso particular) Suponha que $\sum a_n$
convirja. Prove que $\lim_{x \to 1^-} \sum_{n} a_n x^n = \sum_n
a_n$. \emph{(\omterm{thm:b2:series:abel}{Transformação de Abel}: com $A_n$ as somas parciais e $A =
\lim A_n$, escreva $\sum a_nx^n = (1 - x)\sum A_n x^n$; depois $\sum
a_nx^n - A = (1-x)\sum (A_n - A)x^n$, separe a soma num $N$
grande.)} Aplicação: $\sum \frac{(-1)^{n-1}}{n} = \ln 2$ e $\sum
\frac{(-1)^n}{2n+1} = \frac\pi4$, redemonstradas a partir da série de potências.
\end{exercise}

\begin{exercise}[$\star$]\label{exo:b2:powerseries:10}
Mostre que $\displaystyle\sum_{n\geq1}\frac{x^n}{n(n+1)} = 1 +
\frac{1-x}{x}\,\ln(1-x)$ para $0 < \abs x < 1$, determine o
raio e verifique que a convergência é normal em
$\intcc{-1}{1}$; verifique que o valor em $x = 1$ previsto pela
\omterm{def:b2:metric:continuity}{continuidade} coincide com a soma telescópica $\sum
\frac{1}{n(n+1)} = 1$.
\end{exercise}

\begin{exercise}[$\star\star$]\label{exo:b2:powerseries:11}
(Desarranjos) Seja $D_n$ o número de permutações de $n$
objetos sem ponto fixo ($D_0 = 1$). Classificar as
permutações de $\{1, \dots, n\}$ por seu conjunto de pontos fixos dá
$n! = \sum_{k=0}^{n}\binom nk D_{n-k}$. Multiplique por
$\frac{x^n}{n!}$, some e reconheça um \omterm{thm:b2:series:fubini}{produto de Cauchy} para obter
a \omterm{ex:b2:powerseries:fibonacci}{função geradora} exponencial
\[
\sum_{n\geq0} D_n\,\frac{x^n}{n!} = \frac{\eu^{-x}}{1-x}
\qquad (\abs x < 1),
\]
e depois leia a forma fechada $\dfrac{D_n}{n!} =
\sum_{k=0}^{n}\dfrac{(-1)^k}{k!}$ e o limite $\dfrac{D_n}{n!}
\to \eu^{-1}$.
\end{exercise}

\begin{exercise}[$\star\star\star$]\label{exo:b2:powerseries:12}
Prove, com a série binomial do
\cref{met:b2:powerseries:ode}, que
\[
\frac{1}{\sqrt{1 - 4x}} = \sum_{n\geq0}\binom{2n}{n}x^n
\qquad \Bigl(\abs x < \frac14\Bigr),
\]
e deduza, elevando ao quadrado (\omterm{thm:b2:series:fubini}{produto de Cauchy} contra
$\frac{1}{1-4x} = \sum 4^nx^n$), a identidade de convolução
\[
\sum_{k=0}^{n}\binom{2k}{k}\binom{2n-2k}{n-k} = 4^n .
\]
\end{exercise}

\section{Problema: Abel, Tauber e a fronteira da convergência}

\begin{problem}\label{pb:b2:powerseries:1}
Dentro do disco de convergência tudo é fácil; todo o
drama das séries de potências acontece \emph{na} fronteira. Este
problema constrói a teoria da fronteira na variável real: o teorema de
Abel em sua forma uniforme, sua recíproca sob a condição de Tauber,
a hierarquia Cesàro--Abel dos métodos de somação
(com o teorema de Frobenius), a integração termo a termo até a
fronteira com constantes clássicas de dividendo e, enfim,
a rigidez das \omterm{def:b2:powerseries:analytic}{funções analíticas} --- o teorema da
identidade.\index{teorema de Abel}\index{teorema de Tauber}
Em todo o problema, $(a_n)$ é uma sequência real, $f(x) = \sum_{n\geq0}
a_nx^n$ e $A_n = a_0 + \dots + a_n$.

\medskip
\noindent\textbf{Parte I --- O teorema de Abel, \omterm{def:b2:funcseq:def}{uniformemente}.} Suponha,
nesta parte, que $\sum a_n$ convirja, e ponha $r_n =
\sum_{k\geq n} a_k$ (de modo que $r_n \to 0$ e $a_n = r_n -
r_{n+1}$).
\begin{enumerate}
  \item Prove, por somação por partes, que, para todos $0 \leq x
        \leq 1$ e $N \leq M$:
        \[
        \Bigl|\sum_{n=N}^{M} a_n x^n\Bigr|
        \leq 2\sup_{n \geq N}\,\abs{r_n} .
        \]
  \item Deduza que $\sum a_nx^n$ converge \emph{\omterm{def:b2:funcseq:def}{uniformemente}} em
        $\intcc{0}{1}$, que sua soma é \omterm{def:b2:metric:continuity}{contínua} aí e
        recupere o limite radial de
        \cref{exo:b2:powerseries:9}: $f(x) \to \sum a_n$ quando $x
        \to 1^-$.
  \item (Teorema de Abel para \omterm{thm:b2:series:fubini}{produtos de Cauchy}) Sejam $\sum a_n = A$,
        $\sum b_n = B$ e suponha que o \omterm{thm:b2:series:fubini}{produto de Cauchy} $\sum
        c_n$, $c_n = \sum_{k} a_kb_{n-k}$, \emph{convirja},
        com soma $C$. Prove que $C = AB$ \emph{(dentro do disco a
        identidade do produto vale pela
        \cref{prop:b2:powerseries:operations}; faça $x \to
        1^-$)}.
  \item Mostre que a hipótese importa: para $a_n = b_n =
        \frac{(-1)^n}{\sqrt{n+1}}$, as duas séries convergem, e no entanto
        $\abs{c_n} \geq \frac{2(n+1)}{n+2} \geq 1$ \emph{(majore
        cada fator $\sqrt{(k+1)(n-k+1)}$ por MA--MG)}: o
        \omterm{thm:b2:series:fubini}{produto de Cauchy} de duas séries convergentes pode divergir.
  \item (Um dividendo do~\cref{exo:b2:powerseries:5}) Mostre que
        $\bigl(\ln(1-x)\bigr)^2 = 2\sum_{n\geq1}
        \frac{H_n}{n+1}x^{n+1}$ em $\intoo{-1}{1}$, verifique que
        $\bigl(\frac{H_n}{n+1}\bigr)_{n\geq1}$ decresce para
        $0$ e conclua com Abel:
        \[
        \sum_{n\geq1} (-1)^{n+1}\,\frac{H_n}{n+1}
        = \frac{(\ln 2)^2}{2} .
        \]
\end{enumerate}

\medskip
\noindent\textbf{Parte II --- A recíproca de Tauber.} Diga que $\sum
a_n$ é \emph{Abel-somável} para $L$ quando $f(x) \to L$ quando $x \to
1^-$.
\begin{enumerate}[resume]
  \item Mostre que $\sum (-1)^n$ é Abel-somável para $\frac12$
        e no entanto divergente: o teorema de Abel não tem recíproca
        incondicional.
  \item (Lema de Cesàro) Se $u_n \to 0$ então $\frac{u_1 + \dots
        + u_N}{N} \to 0$ \emph{(separe a soma num $m$ fixo)}.
  \item Suponha agora $n\,a_n \to 0$ e $f(x) \to L$. Com $x_N =
        1 - \frac1N$, prove as duas estimativas
        \[
        \Bigl|\sum_{n=0}^{N} a_n\bigl(1 - x_N^n\bigr)\Bigr|
        \leq \frac{1}{N}\sum_{n=1}^{N} n\,\abs{a_n},
        \qquad
        \Bigl|\sum_{n>N} a_n x_N^n\Bigr|
        \leq \sup_{n>N}\bigl(n\abs{a_n}\bigr)
        \]
        \emph{(para a primeira, $1 - x^n \leq n(1-x)$; para a
        segunda, $\abs{a_n} \leq \frac{1}{N}\sup_{m>N}
        m\abs{a_m}$ e $\sum x_N^n \leq N$)}.
  \item Conclua o \emph{teorema de Tauber}: se $n\,a_n \to 0$ e
        $\sum a_n$ é Abel-somável para $L$, então $\sum a_n$
        converge para $L$.
  \item (O tauberiano fácil para coeficientes positivos) Se $a_n
        \geq 0$ e $f$ é limitada em $\intco{0}{1}$, mostre que
        $\sum a_n$ converge e $\sum a_n = \lim_{x\to1^-}
        f(x)$ \emph{(majore $\sum_{n\leq N}a_nx^n \leq f(x)$ e
        faça $x \to 1^-$, depois use Abel)}.
\end{enumerate}

\medskip
\noindent\textbf{Parte III --- Médias de Cesàro e o teorema de
Frobenius.} Diga que $\sum a_n$ é \emph{Cesàro-somável} para $L$ quando
$\sigma_N = \frac{A_0 + \dots + A_{N-1}}{N} \to L$.
\begin{enumerate}[resume]
  \item Mostre que uma série convergente é Cesàro-somável para
        sua soma \emph{(questão 7 aplicada a $A_n - L$)}.
  \item Calcule o valor de Cesàro de $\sum(-1)^n$ e verifique que ele
        coincide com o valor de Abel $\frac12$ da questão 6.
  \item Com $S_n = A_0 + \dots + A_n = (n+1)\,\sigma_{n+1}$,
        prove as duas identidades, para $0 \leq x < 1$:
        \[
        f(x) = (1-x)^2\sum_{n\geq0}(n+1)\,\sigma_{n+1}x^n,
        \qquad
        (1-x)^2\sum_{n\geq0}(n+1)x^n = 1 .
        \]
  \item (Frobenius) Deduza: se $\sigma_N \to L$ então $f(x) \to
        L$ quando $x \to 1^-$ --- Cesàro-somável implica
        Abel-somável, com o mesmo valor \emph{(subtraia as duas
        identidades e separe a soma num $N$ grande, como na
        \cref{exo:b2:powerseries:9})}.
  \item Mostre que a hierarquia
        \[
        \text{convergente} \;\Longrightarrow\;
        \text{Cesàro-somável} \;\Longrightarrow\;
        \text{Abel-somável}
        \]
        é estrita nas duas setas: a questão 6 para a primeira;
        para a segunda, mostre que $\sum(-1)^n(n+1)$ é
        Abel-somável para $\frac14$ (calcule $f$) mas não
        Cesàro-somável (calcule $\sigma_N$ separadamente para
        $N$ par e ímpar).
\end{enumerate}

\medskip
\noindent\textbf{Parte IV --- Integrando até a fronteira.}
\begin{enumerate}[resume]
  \item Suponha que $\sum a_nx^n$ convirja em $\intco{0}{1}$ e que
        $\sum \frac{a_n}{n+1}$ convirja. Prove que a
        \omterm{def:b2:integration:improper}{integral imprópria} $\int_0^1 f$ existe e
        \[
        \int_0^1 \Bigl(\sum_{n\geq0} a_nx^n\Bigr)\dd x
        = \sum_{n\geq0}\frac{a_n}{n+1}
        \]
        \emph{(a primitiva $F(x) = \sum\frac{a_n}{n+1}x^{n+1}$
        é \omterm{def:b2:metric:continuity}{contínua} em $1$ pela Parte I)}.
  \item Seja $\eta = \sum_{n\geq1}\frac{(-1)^{n-1}}{n^2}$. Mostre que
        $\int_0^1 \frac{\ln(1+x)}{x}\dd x = \eta$ e, separando índices pares e ímpares na
        série \omterm{def:b2:series:def}{absolutamente} convergente $\sum \frac1{n^2}$, que $\eta =
        \frac12\sum_{n\geq1}\frac{1}{n^2}$. (O problema de fim de semana do
        capítulo de Fourier avalia
        $\sum\frac1{n^2} = \frac{\pi^2}{6}$.)
  \item Prove
        \[
        \sum_{n\geq0}\frac{(-1)^n}{3n+1}
        = \int_0^1\frac{\dd x}{1+x^3}
        = \frac13\Bigl(\ln 2 + \frac{\pi}{\sqrt3}\Bigr)
        \]
        \emph{(a série converge por Leibniz; integre a
        série geométrica $\sum(-1)^nx^{3n}$ com a questão 16;
        depois frações parciais: $\frac{1}{1+x^3} =
        \frac{1/3}{1+x} + \frac{(2-x)/3}{x^2-x+1}$)}.
  \item Da série binomial de $(1-t)^{-1/2}$
        (\cref{exo:b2:powerseries:12}) deduza
        \[
        \arcsin x = \sum_{n\geq0}
        \frac{\binom{2n}{n}}{4^n(2n+1)}\,x^{2n+1}
        \quad(\abs x < 1),
        \qquad\text{então}\qquad
        \sum_{n\geq0}\frac{\binom{2n}{n}}{4^n(2n+1)}
        = \frac\pi2 ,
        \]
        justificando o valor na fronteira pela convergência
        \emph{normal} em $\intcc{-1}{1}$ (use
        $\binom{2n}n4^{-n} \sim \frac{1}{\sqrt{\pi n}}$,
        \cref{ex:b2:comparison:centralbinomial}) --- aqui nem
        Abel é necessário.
  \item (Catalan na fronteira) Mostre que $\sum C_n 4^{-n} =
        2$: a série de Catalan do
        \cref{ex:b2:powerseries:catalan} converge \emph{no}
        seu raio $\frac14$ (assintótica de
        \cref{exo:b2:powerseries:8}), sua soma é \omterm{def:b2:metric:continuity}{contínua} em
        $\intcc{0}{\frac14}$, e a forma fechada tem limite $2$
        aí.
\end{enumerate}

\medskip
\noindent\textbf{Parte V --- Rigidez: o teorema da identidade.}
\begin{enumerate}[resume]
  \item (Zeros isolados) Seja $f = \sum a_nx^n$ de raio $R >
        0$ com nem todos os $a_n = 0$ nulos; seja $m$ o menor índice
        com $a_m \neq 0$. Mostre que $f(x) = x^m g(x)$ com $g$ uma
        série de potências de raio $R$, $g(0) = a_m \neq 0$, e
        deduza que $f$ não tem zero em alguma
        vizinhança perfurada de $0$.
  \item (Teorema da identidade) Sejam $f, h$ somas de séries de potências
        perto de $0$ e $(x_k)$ uma sequência de pontos \emph{não nulos}
        com $x_k \to 0$ e $f(x_k) = h(x_k)$. Prove
        que $f$ e $h$ têm os mesmos coeficientes, logo
        coincidem perto de $0$.
  \item Determine \emph{todas} as funções $f$ \omterm{def:b2:powerseries:analytic}{analíticas} perto de $0$ com
        \[
        f\Bigl(\frac1k\Bigr) = \frac{k^2}{k^2+1}
        \qquad\text{para todo inteiro grande } k .
        \]
  \item Mostre que uma \omterm{def:b2:powerseries:analytic}{função analítica} num intervalo \omterm{def:b2:metric:topology}{aberto} $I$
        que se anula num subintervalo se anula identicamente em
        $I$ \emph{(o conjunto dos pontos em torno dos quais $f$ se anula
        identicamente é \omterm{def:b2:metric:topology}{aberto} e, pelo teorema da identidade
        aplicado em pontos de acumulação, fechado em $I$)}.
        Conclua que nenhuma \omterm{def:b2:powerseries:analytic}{função analítica} não nula em $\R$ tem
        suporte \omterm{def:b2:metric:compact}{compacto} --- ao passo que existem funções bolha
        $C^\infty$ (o~\cref{exo:b2:powerseries:7} fornece o
        bloco de construção): a analiticidade é rígida, a suavidade é
        frouxa.
  \item Síntese. Uma frase para cada: (i) o que o teorema de Abel
        acrescenta ao pacote de \omterm{def:b2:funcseq:series}{convergência normal} do
        \cref{lem:b2:powerseries:abel}; (ii) as hipóteses exatas
        sob as quais a recíproca vale (Tauber) e
        o degrau intermediário (Frobenius); (iii) uma constante
        de fronteira da Parte IV que você agora saberia deduzir para um
        amigo em duas linhas; (iv) onde as médias de Cesàro
        reaparecerão neste livro, para séries de tipo bem
        diferente.
\end{enumerate}
\end{problem}
